
\documentclass{article}
\usepackage{colortbl}
\usepackage{makecell}
\usepackage{multirow}
\usepackage{supertabular}

\begin{document}

\newcounter{utterance}

\centering \large Interaction Transcript for game `matchit\_ascii', experiment `similar\_grid\_1', episode 5 with NEW\_Reflect\_Agent.
\vspace{24pt}

{ \footnotesize  \setcounter{utterance}{1}
\setlength{\tabcolsep}{0pt}
\begin{supertabular}{c@{$\;$}|p{.15\linewidth}@{}p{.15\linewidth}p{.15\linewidth}p{.15\linewidth}p{.15\linewidth}p{.15\linewidth}}
    \# & $\;$A & \multicolumn{4}{c}{Game Master} & $\;\:$B\\
    \hline

    \theutterance \stepcounter{utterance}  
    & & \multicolumn{4}{p{0.6\linewidth}}{
        \cellcolor[rgb]{0.9,0.9,0.9}{
            \makecell[{{p{\linewidth}}}]{
                \texttt{\tiny{[P1$\langle$GM]}}
                \texttt{You are participating in a collaborative guessing game.} \\
\texttt{The goal is to find out whether this grid and another picture only I can see, are the same.} \\
\texttt{Please describe your grid first. Then, I will provide my description and we can ask each other questions about the grids to figure out whether they are the same. Now start your short grid description with "DESCRIPTION:" followed by the description. Do not add anything else. This is your grid:} \\
\texttt{X ▢ ▢ ▢ X} \\
\texttt{X X ▢ ▢ X} \\
\texttt{X ▢ X ▢ X} \\
\texttt{X ▢ ▢ X X} \\
\texttt{X ▢ ▢ ▢ X} \\
            }
        }
    }
    & \\ \\

    \theutterance \stepcounter{utterance}  
    & \multicolumn{4}{p{0.45\linewidth}}{
        \cellcolor[rgb]{0.95,0.95,0.95}{
            \makecell[{{p{\linewidth}}}]{
                \texttt{\tiny{[P1$|$internal]}}
                \texttt{Reflections from previous episodes:} \\
\texttt{{-}} \\
\texttt{ Failed. I gave a clear description of my grid pattern, but the user did not provide their description or engage in questioning, ending the episode prematurely. Next time, I would follow my initial description by immediately prompting for their description and proposing specific comparative questions to drive the collaboration forward.} \\
\texttt{{-}} \\
\texttt{ Failed. I described the grid as having a vertical line of X's in the third column, but I didn't specify the exact pattern of the first row (all X's). This could cause mismatch if the other grid differed in the top row. Next time, I would give a complete description including any full rows or columns, not just the central feature.} \\
\texttt{{-}} \\
\texttt{ Failed. I only gave my initial description but did not engage in questioning or compare with the partner's description. Next time, after providing my description, I must actively ask questions about the other grid and answer any questions about mine to determine if they match.} \\
            }
        }
    }
    & & \\ \\

    \theutterance \stepcounter{utterance}  
    & \multicolumn{4}{p{0.45\linewidth}}{
        \cellcolor[rgb]{0.95,0.95,0.95}{
            \makecell[{{p{\linewidth}}}]{
                \texttt{\tiny{[P1$|$internal]}}
                \texttt{{[}attempt 1 – valid{]}} \\
\texttt{DESCRIPTION: The grid is 5x5. The first column is all X's. The last column is all X's. The first row has X's in the first and last positions, with ▢ in the middle three. The second row has X's in the first two and last positions, with ▢ in the middle two. The third row has X's in the first, third, and last positions, with ▢ in the second and fourth. The fourth row has X's in the first, fourth, and fifth positions, with ▢ in the second and third. The fifth row has X's in the first and last positions, with ▢ in the middle three.} \\
            }
        }
    }
    & & \\ \\

    \theutterance \stepcounter{utterance}  
    & \multicolumn{4}{p{0.6\linewidth}}{
        \cellcolor[rgb]{0.9,0.9,0.9}{
            \makecell[{{p{\linewidth}}}]{
                \texttt{\tiny{[P1$\rangle$GM]}}
                \texttt{DESCRIPTION: The grid is 5x5. The first column is all X's. The last column is all X's. The first row has X's in the first and last positions, with ▢ in the middle three. The second row has X's in the first two and last positions, with ▢ in the middle two. The third row has X's in the first, third, and last positions, with ▢ in the second and fourth. The fourth row has X's in the first, fourth, and fifth positions, with ▢ in the second and third. The fifth row has X's in the first and last positions, with ▢ in the middle three.} \\
            }
        }
    }
    & & \\ \\

    \theutterance \stepcounter{utterance}  
    & & & \multicolumn{2}{p{0.3\linewidth}}{
        \cellcolor[rgb]{0.9,0.9,0.9}{
            \makecell[{{p{\linewidth}}}]{
                \texttt{\tiny{[GM$|$GM]}}
                \texttt{continue} \\
            }
        }
    }
    & & \\ \\

    \theutterance \stepcounter{utterance}  
    & & \multicolumn{4}{p{0.6\linewidth}}{
        \cellcolor[rgb]{0.9,0.9,0.9}{
            \makecell[{{p{\linewidth}}}]{
                \texttt{\tiny{[GM$\rangle$P2]}}
                \texttt{You are participating in a collaborative guessing game.} \\
\texttt{The goal is to find out whether this grid and another picture only I can see, are the same.} \\
\texttt{Please describe your grid first. Then, I will provide my description and we can ask each other questions about the grids to figure out whether they are the same. Now start your short grid description with "DESCRIPTION:" followed by the description. Do not add anything else. This is your grid:} \\
\texttt{X ▢ ▢ ▢ X} \\
\texttt{X ▢ ▢ X X} \\
\texttt{X ▢ X ▢ X} \\
\texttt{X X ▢ ▢ X} \\
\texttt{X ▢ ▢ ▢ X} \\
            }
        }
    }
    & \\ \\

    \theutterance \stepcounter{utterance}  
    & & & \multicolumn{4}{p{0.45\linewidth}}{
        \cellcolor[rgb]{0.95,0.95,0.95}{
            \makecell[{{p{\linewidth}}}]{
                \texttt{\tiny{[P2$|$internal]}}
                \texttt{Reflections from previous episodes:} \\
\texttt{{-}} \\
\texttt{ FAIL. My description was inaccurate: I described a plus sign, but my grid has an extra X at (4,2) and (4,4), making the shape more like a cross with a thicker base or a small "T" shape, not a symmetric plus. This mismatch likely caused confusion. Next time, I would count X's precisely and describe exact coordinates of deviations from the simple plus, e.g., "plus sign with two additional X's symmetrically below the center row."} \\
\texttt{{-}} \\
\texttt{ Failed. My description incorrectly stated the first row was all X's, but my grid actually had only the third column as X in rows 1{-}4, with a full bottom row of X's. This major error misled the comparison from the start. Next time, I must describe the grid accurately and systematically, row by row, to avoid such a critical mistake.} \\
\texttt{{-}} \\
\texttt{ FAIL. My description incorrectly placed X's in the first and fifth positions of the first row and second/fourth of the second row, which do not exist in my actual grid. I described a pattern from my memory/assumption instead of directly reading the provided grid. This introduced major discrepancies from the start.} \\
\\ 
\texttt{Next time: Read the grid precisely row by row for the description. For example, correctly state: "Top row: empty, empty, X, empty, empty. Second row: same. Third row: same. Fourth row: empty, X, empty, X, empty. Bottom row: X, empty, empty, empty, X."} \\
            }
        }
    }
     \\ \\

    \theutterance \stepcounter{utterance}  
    & & & \multicolumn{4}{p{0.45\linewidth}}{
        \cellcolor[rgb]{0.95,0.95,0.95}{
            \makecell[{{p{\linewidth}}}]{
                \texttt{\tiny{[P2$|$internal]}}
                \texttt{{[}attempt 1 – valid{]}} \\
\texttt{DESCRIPTION: Top row: X, empty, empty, empty, X. Second row: X, empty, empty, X, X. Third row: X, empty, X, empty, X. Fourth row: X, X, empty, empty, X. Bottom row: X, empty, empty, empty, X.} \\
            }
        }
    }
     \\ \\

    \theutterance \stepcounter{utterance}  
    & & & \multicolumn{4}{p{0.6\linewidth}}{
        \cellcolor[rgb]{0.9,0.9,0.9}{
            \makecell[{{p{\linewidth}}}]{
                \texttt{\tiny{[GM$\langle$P2]}}
                \texttt{DESCRIPTION: Top row: X, empty, empty, empty, X. Second row: X, empty, empty, X, X. Third row: X, empty, X, empty, X. Fourth row: X, X, empty, empty, X. Bottom row: X, empty, empty, empty, X.} \\
            }
        }
    }
     \\ \\

    \theutterance \stepcounter{utterance}  
    & & & \multicolumn{2}{p{0.3\linewidth}}{
        \cellcolor[rgb]{0.9,0.9,0.9}{
            \makecell[{{p{\linewidth}}}]{
                \texttt{\tiny{[GM$|$GM]}}
                \texttt{continue} \\
            }
        }
    }
    & & \\ \\

    \theutterance \stepcounter{utterance}  
    & & \multicolumn{4}{p{0.6\linewidth}}{
        \cellcolor[rgb]{0.9,0.9,0.9}{
            \makecell[{{p{\linewidth}}}]{
                \texttt{\tiny{[GM$\rangle$P2]}}
                \texttt{This is my DESCRIPTION: The grid is 5x5. The first column is all X's. The last column is all X's. The first row has X's in the first and last positions, with ▢ in the middle three. The second row has X's in the first two and last positions, with ▢ in the middle two. The third row has X's in the first, third, and last positions, with ▢ in the second and fourth. The fourth row has X's in the first, fourth, and fifth positions, with ▢ in the second and third. The fifth row has X's in the first and last positions, with ▢ in the middle three.} \\
\texttt{Now ask a question in order to find out new aspects of my grid that may be different to your grid. Start with "QUESTION:" and do not add anything else.} \\
            }
        }
    }
    & \\ \\

    \theutterance \stepcounter{utterance}  
    & & & \multicolumn{4}{p{0.6\linewidth}}{
        \cellcolor[rgb]{0.9,0.9,0.9}{
            \makecell[{{p{\linewidth}}}]{
                \texttt{\tiny{[GM$\langle$P2]}}
                \texttt{DESCRIPTION: The grid is 5x5. The first column is all X's. The last column is all X's. The first row has X's in the first and last positions, with ▢ in the middle three. The second row has X's in the first two and last positions, with ▢ in the middle two. The third row has X's in the first, third, and last positions, with ▢ in the second and fourth. The fourth row has X's in the first, fourth, and fifth positions, with ▢ in the second and third. The fifth row has X's in the first and last positions, with ▢ in the middle three.  } \\
\texttt{QUESTION: Is the second row's third cell an X or a ▢?} \\
            }
        }
    }
     \\ \\

    \theutterance \stepcounter{utterance}  
    & & & \multicolumn{2}{p{0.3\linewidth}}{
        \cellcolor[rgb]{0.9,0.9,0.9}{
            \makecell[{{p{\linewidth}}}]{
                \texttt{\tiny{[GM$|$GM]}}
                \texttt{abort, first word: DESCRIPTION:, but expected QUESTION:} \\
            }
        }
    }
    & & \\ \\

\end{supertabular}
}

\end{document}
